\documentclass[../../main.tex]{subfiles}
\begin{document}

\begin{center}
  {\fontsize{48pt}{58pt}\selectfont \textbf{1996}\normalsize\textbf{年度}}
\end{center}

\begin{table}[htb]
  \centering
  \caption{出題テーマと難易度}
  \begin{tabular}{|c|p{9cm}|c|}
    \hline
    \textbf{問題番号} & \textbf{テーマ} & \textbf{難易度} \\
    \hline
    第1問 & 円上の点から定まる写像と回転数 & \\
    \hline
    第2問 & 四面体の平面切断の最大面積 & \\
    \hline
    第3問 & 一次変換の3回合成 & \\
    \hline
    第4問 & 床関数を使う数列の収束 & \\
    \hline
    第5問 & 正三角形内の点の反復操作と期待値 & \\
    \hline
    第6問 & 速度と燃料消費量の最適化 & \\
    \hline
  \end{tabular}
\end{table}

\begin{table}[htb]
  \centering
  \caption{解答}
  \renewcommand{\arraystretch}{1.4}
  \begin{tabular}{|c|c|p{9cm}|}
    \hline
    \textbf{問題番号} & \textbf{小問} & \textbf{答え} \\
    \hline
    第1問 & (1) & $|\overrightarrow{OY}|=1$（証明問題） \\
    \cline{2-3}
    & (2) & $X=A$または$A$の対蹠点 \\
    \cline{2-3}
    & (3) & $X$が1周すると$Y$は2周する（証明問題） \\
    \hline
    第2問 & -- & 最大値$\dfrac{3\sqrt{2}-4}{2}$（$t=\dfrac{1+\sqrt{2}}{2}$） \\
    \hline
    第3問 & -- &  \\
    \hline
    第4問 & (1) & $k=8$では$0,2,3,3,\ldots$，$k=9$では$0,3,4,4,\ldots$ \\
    \cline{2-3}
    & (2) & $a_n\le\dfrac{k-1}{2}$かつ$a_n\le a_{n+1}$（証明問題） \\
    \cline{2-3}
    & (3) & 以後一定．値は$k$が奇数なら$(k-1)/2$，偶数なら$k/2-1$ \\
    \hline
    第5問 & (1) & $E_1=\dfrac{1+|\overrightarrow{OP_0}|^2}{4}$ \\
    \cline{2-3}
    & (2) & $\overrightarrow{OP_n}=2^{-n}\overrightarrow{OP_0}+\sum_{i=1}^{n}2^{i-n-1}\overrightarrow{OX_i}$ \\
    \cline{2-3}
    & (3) & $E_n=\dfrac{1}{3}\left(1-4^{-n}\right)$ \\
    \hline
    第6問 & -- & 初めに$100(e^{ek}-1)\,\mathrm{kg}$積み，時速$1/k\,\mathrm{km}$で走る \\
    \hline
  \end{tabular}
\end{table}

\paragraph{試験形態}
6問．試験時間は要確認．

\end{document}
